\section{Chain of Custody, Cryptography and Verifiability (RQ3)}
\label{sec:custody_crypto}

To resolve RQ3, this section examines how digital forensics reconciles probabilistic neural inference with statutory custody, deterministic replay, examiner impartiality, and multi-jurisdictional data protection mandates.

\subsection{Normative Alignment and Cryptographic Provenance}
End-to-end digital evidence integrity spans three complementary international standards: ISO/IEC 27037~\cite{iso27037_2012} governs identification, collection, acquisition, and preservation; ISO/IEC 27041~\cite{iso27041_2015} defines assurance criteria for validating investigative methods; and ISO/IEC 27042~\cite{iso27042_2015} establishes procedural guidelines for the analysis and interpretation phase where language models operate, complemented by NIST SP 800-86~\cite{nist_sp800_86}. In conventional forensics, integrity relies on computing cryptographic digests (SHA-256 or BLAKE3) immediately upon disk imaging or memory capture. Any bit-level alteration invalidates the evidentiary baseline.

Introducing generative inference into the analysis phase threatens this continuous audit trail whenever model queries execute as ephemeral, unlogged interactions. Defensible forensic engineering therefore requires treating every inference call as an auditable, cryptographically sealed transaction~\cite{sharma2025_forensicllm, chernyshev2026_large, tiwari2026_xforensicnet}. Specifically, the system prompt, serialized evidence slice, model checkpoint digest, quantization specification, decoding configuration, and raw generated output tokens must be hashed and chained into an append-only provenance ledger signed by the forensic workstation. Should opposing counsel challenge an automated triage note, an independent auditor can verify that the exact evidence slice and model artifact existed unaltered at the recorded timestamp.

\subsection{Inference Determinism, Daubert, and FRE 702(d)}
In common-law jurisdictions governed by the Daubert standard~\cite{daubert1993} and Federal Rule of Evidence 702 (as amended in 2023)~\cite{fre702_2023}, scientific expert methodology is screened for empirical testability, peer review, known or potential error rates, general acceptance (echoing Frye~\cite{frye1923}), and reliable application of principles to case-specific facts under Rule 702(d). Reproducibility failures have already surfaced in litigation: as documented by Trivedi~\cite{trivedi2026_when}, in \textit{Matter of Weber} (2024) a surrogate's court submitted the same Microsoft Copilot query across three computers and obtained three divergent figures, while the testifying expert could not recall the exact prompt used.

Practitioners frequently assume that setting greedy decoding temperature to zero ($T=0.0$) guarantees exact bitwise replay. Hardware execution physics refutes this assumption~\cite{michelet2024_chatgpt, trivedi2026_when, scanlon2026_investigation}. Modern GPU tensor libraries (such as NVIDIA cuBLAS and cuDNN) execute parallel warp reductions whose floating-point summation order varies dynamically with thread scheduling and batch composition. Because IEEE 754 floating-point addition in FP16, BF16, and FP32 is non-associative, microscopic rounding variations ($\approx 10^{-6}$) alter least significant mantissa bits in logit vectors. Whenever two candidate vocabulary tokens exhibit near-identical logits, thread-scheduling jitter flips the greedy $\arg\max$ selection; once an alternate token enters the autoregressive context window, subsequent text diverges rapidly. Deterministic replay remains unattainable on standard batched GPU pipelines.

To establish auditable execution verifiability under ISO/IEC 27041 and FRE 702, forensic architectures require a complete cryptographic execution manifest $\mathcal{M}$:
\begin{align}
\label{eq:forensic_manifest}
\mathcal{M} = \langle &\mathcal{H}_{\text{weights}}, \mathcal{H}_{\text{vocab}}, \mathcal{H}_{\text{tmpl}}, \mathcal{H}_{\text{prompt}}, \mathcal{H}_{\text{evidence}}, \nonumber \\
&q_{\text{quant}}, \theta_{\text{dec}}, \sigma_{\text{seed}}, \phi_{\text{runtime}}, \kappa_{\text{det}} \rangle
\end{align}
Here, $\mathcal{H}_{\text{weights}}$, $\mathcal{H}_{\text{vocab}}$, $\mathcal{H}_{\text{tmpl}}$, $\mathcal{H}_{\text{prompt}}$, and $\mathcal{H}_{\text{evidence}}$ denote immutable SHA-256 digests of the model checkpoint, tokenizer vocabulary, chat template, system/user prompt, and input evidence chunk; $q_{\text{quant}}$ records the exact quantization format (e.g., GGUF $Q4\_K\_M$ or FP16); $\theta_{\text{dec}} = (T, \text{top-}p, \text{top-}k, \text{rep\_pen})$ fixes all decoding hyperparameters; $\sigma_{\text{seed}}$ pins the pseudo-random seed; $\phi_{\text{runtime}}$ logs exact engine, compiler, and driver versions; and $\kappa_{\text{det}}$ enforces runtime-specific determinism controls (e.g., \texttt{CUBLAS\_WORKSPACE\_CONFIG=:4096:8} and \texttt{torch.use\_deterministic\_algorithms(True)} for PyTorch/vLLM backends, or fixed single-sequence batch size, deterministic thread affinity, and \texttt{--temp 0 --seed} pinning for \texttt{llama.cpp}/GGML engines). Single-threaded CPU execution eliminates GPU warp non-determinism entirely, yet its low token throughput renders it operationally impractical for bulk multi-terabyte triage; instead, laboratories should combine batched local GPU triage with targeted deterministic CPU verification strictly on short, dispositive evidence excerpts selected for courtroom presentation.

Table~\ref{tab:daubert_mapping} maps Daubert and FRE 702 admissibility pillars directly against LLM technical constraints and concrete engineering safeguards.

\begin{table*}[!t]
\caption{Admissibility Assessment: Daubert Standard and FRE 702(d) Criteria vs. Large Language Model Realities in Digital Forensics}
\label{tab:daubert_mapping}
\centering
\scriptsize
\begin{tabular}{p{2.6cm}p{4.2cm}p{4.4cm}p{4.4cm}}
\toprule
\textbf{Legal Criterion (Daubert / FRE 702)} & \textbf{Forensic Baseline Requirement} & \textbf{LLM Inherent Vulnerability} & \textbf{Forensic Engineering Mitigation} \\
\midrule
\textbf{Empirical Testability} & Methodology must be falsifiable under ISO/IEC 27041 and standardized CFTT-style benchmarks. & Emergent latent reasoning evades closed-form verification; behavior shifts across model updates. & Decouple triage from fact verification; corroborate every claim via deterministic regex, hex offsets, and hashes. \\
\addlinespace
\textbf{Known or Potential Error Rate} & Quantifiable error bounds and reproducible false-positive/negative metrics. & Error rates are prompt-dependent rather than fixed method properties; GPU non-associativity induces drift. & Record full manifest $\mathcal{M}$; run fixed-batch local inference with targeted deterministic CPU replay on key artifacts. \\
\addlinespace
\textbf{Peer Review and Publication} & Underlying algorithms, codebases, and validation recipes must undergo public scrutiny. & Closed commercial APIs (GPT-4, Claude, Gemini) conceal weights, training mixtures, and silent updates. & Mandate archival open-weights local models (e.g., ForensicLLM, Llama variants) with pinned SHA-256 checkpoints. \\
\addlinespace
\textbf{Reliable Application (FRE 702(d)) \& General Acceptance} & Reliable case-specific application and scientific consensus among practitioners. & Uncritical acceptance of unsupported premises, hallucinations, and prompt-injection exposure. & Restrict LLMs to assistive copilot triage; mandate premise-audit prompting and sworn human verification. \\
\bottomrule
\end{tabular}
\end{table*}

\subsection{Premise Uncriticality, Sycophancy, and Non-Delegable Accountability}
Beyond floating-point drift, recent empirical jurisprudence exposes a nuanced cognitive vulnerability in AI-assisted forensic reasoning. Forensic examiners carry strict statutory duties of impartiality, evaluating both inculpatory and exculpatory hypotheses objectively. While early commentary warned that conversational models exhibit straightforward sycophancy by echoing leading investigator prompts, rigorous empirical testing reveals a sharper failure mechanism. Trivedi~\cite{trivedi2026_when} tested six dated model snapshots across five forensic interpretation items under three examiner framings (2,700 total trials; 2,160 analyzed, inter-rater Gwet's $\text{AC1}=0.855$). Contrary to naive sycophancy assumptions, the direct assertion hypothesis was not supported: models asserted the incorrect proposition in 36.1\% of neutral trials compared to 19.6\% when an examiner explicitly asserted it. However, the study uncovered a far more dangerous epistemic flaw: on an item whose conclusion rested on an unsupported premise, five of the six evaluated models failed to contest the premise in 0 of 150 trials (while only one model contested it in 80\% of trials). Because models verify surface conclusions without auditing underlying premises, error rates become prompt-dependent rather than intrinsic properties of a scientific method, rendering unverified LLM outputs directly vulnerable to exclusion under FRE 702(d)~\cite{trivedi2026_when}.

Technical countermeasures mitigate but do not eliminate this epistemic risk. Human-supervised extraction workflows~\cite{ko2026_humansupervised} require analysts to verify every extracted identifier against raw source documents before downstream evidentiary use, while offline neuro-symbolic frameworks such as X-ForensicNet~\cite{tiwari2026_xforensicnet} constrain neural deductions against deterministic Windows forensic rules.

Ultimately, judicial doctrine establishes a hard legal boundary: personal, non-delegable examiner accountability. Across both common-law adversarial courts and continental civil-law jurisdictions, forensic expert testimony remains strictly an \textit{in personam} legal duty borne by a sworn human practitioner subject to civil and criminal perjury liability. In continental criminal procedure, exemplified by Article~163 of the Portuguese Code of Penal Procedure (\textit{C{\'o}digo de Processo Penal}, CPP)~\cite{portugal_cpp}, technical and scientific expert findings enjoy a statutory presumption removing them from unrestricted free judicial evaluation (n.{\textordmasculine}~1); crucially, unlike pre-admission Daubert gatekeeping in U.S.\ federal courts, Article~163(2) requires the judge to state explicit scientific grounds whenever diverging from the expert opinion, while opaque or unverifiable methods expose the report to evidentiary challenge, secondary expert examination, or procedural nullity. An artificial intelligence system holds no legal personality, cannot swear an oath, and cannot undergo cross-examination. Adopting an unverified generative inference into a formal forensic report constitutes examiner negligence. Language models therefore remain non-witness assistive instruments subordinate to human accountability.

\subsection{European Data Protection, EU AI Act, and Air-Gapped Local Deployment}
Seized digital evidence in European investigations intersects three distinct statutory regimes that academic prototypes routinely conflate. First, in corporate incident response, internal audits, and civil litigation, evidentiary processing is governed by Regulation (EU) 2016/679 (General Data Protection Regulation, GDPR)~\cite{gdpr2016}, which enforces purpose limitation (Article~5), heightened safeguards for special categories of personal data frequently present on seized endpoints (Article~9), and strict international transfer controls (Articles~44--49). Second, in criminal investigations conducted by police and judicial authorities, GDPR Article~2(2)(d) explicitly excludes criminal law enforcement; instead, processing is governed by the European Union Law Enforcement Directive, Directive (EU) 2016/680 (LED)~\cite{eu_led_2016}, transposed into Portuguese law by Law n.{\textordmasculine}~59/2019~\cite{portugal_lei59_2019}. Within national criminal procedure, Article~86 of the Portuguese CPP (investigative secrecy, \textit{segredo de justi{\c{c}}a})~\cite{portugal_cpp} and Article~16 of Law n.{\textordmasculine}~109/2009 (Portuguese Cybercrime Law)~\cite{leicibercrime2009} mandate strict confidentiality and preservation of seized computer data. Third, Regulation (EU) 2024/1689 (the EU Artificial Intelligence Act)~\cite{eu_ai_act_2024}, Annex~III, Point~6(d), classifies AI systems intended for use by law enforcement authorities to evaluate the reliability of evidence in criminal investigations or prosecutions as \textbf{high-risk AI systems}, triggering statutory obligations for automatic event logging (Article~12), technical documentation, risk management, accuracy and cybersecurity robustness (Article~15), and effective human oversight (Article~14).

Our orthogonal corpus audit exposes an acute compliance gap between academic experimentation and these statutory mandates. Across the 34 empirical LLM-for-DFIR tool studies, 67.6\% ($n=23$) rely on commercial cloud APIs or omit explicit local air-gapped execution guarantees, whereas only 23.5\% ($n=8$) design or deploy dedicated local, offline, or lightweight open-weights models (with another 8.8\%, $n=3$, including open-weights baselines in multi-model benchmarks). Transmitting unredacted disk artifacts, memory strings, or seized communications to external commercial inference endpoints breaches GDPR/LED cross-border transfer rules, violates statutory judicial secrecy under CPP Article~86 and Law n.{\textordmasculine}~109/2009, severs physical custody accounting under ISO/IEC 27037, and fails EU AI Act high-risk logging and auditability controls.

Reconciling generative assistance with statutory admissibility mandates air-gapped on-premise deployment coupled with hardware write-blocking under ISO/IEC 27037 and analytical isolation under ISO/IEC 27042. Primary evidentiary containers (RAW, E01, or AFF4) must mount exclusively read-only with timestamp preservation (\texttt{mount -o ro,loop,noatime}), while local inference engines (such as offline \texttt{llama.cpp} or air-gapped \texttt{vLLM} instances) execute inside segregated volatile memory scratch partitions so KV-cache allocations and temporary token buffers never contaminate primary evidence media.

Hardware feasibility for regional police laboratories depends on realistic memory budgeting. While 70B-parameter models demand datacenter accelerators (80~GB+ VRAM), 4-bit quantized Small Language Models (SLMs) spanning 7B to 8B parameters (e.g., Llama-3-8B or Mistral-7B in GGUF $Q4\_K\_M$) require approximately 4.5--5.0~GB of VRAM for model weights alone (roughly 5.8~GB total under short 4k--8k token contexts)~\cite{sharma2025_forensicllm, michelet2025_finetuning}. Crucially, however, operational forensic triage over large log slices induces substantial \textbf{KV-cache memory growth}: extending context windows to 32k--128k tokens pushes total VRAM consumption to 12--24~GB depending on Grouped-Query Attention (GQA) configuration and KV-cache precision. Budget-constrained forensic workstations (e.g., 16--24~GB consumer GPUs or unified-memory Apple Silicon hosts) must therefore pair 4-bit weight quantization with KV-cache quantization (Q8\_0) and deterministic fixed-size evidence chunking to preserve both air-gapped privacy and operational stability.

\begin{figure*}[!t]
\centering
\begin{tikzpicture}[
    node distance=0.6cm and 0.32cm,
    stage/.style={draw=blue!60!black, fill=blue!5, rounded corners=2pt, font=\scriptsize\bfseries, align=center, inner sep=3pt, text width=2.45cm, minimum height=1.15cm},
    control/.style={draw=red!60!black, fill=red!5, dashed, rounded corners=2pt, font=\tiny, align=center, inner sep=2pt, text width=2.45cm},
    arr/.style={-{Latex[length=2.2mm]}, thick, draw=blue!60!black}
]

% Nodes
\node[stage] (evid) {1. Evidence Image\\(E01, RAW, AFF4)\\Bitstream Capture};
\node[stage, right=0.32cm of evid] (block) {2. Ingestion Isolation\\Hardware Write-Block\\Read-Only (\texttt{noatime})};
\node[stage, right=0.32cm of block] (parse) {3. Deterministic Parse\\Plaso / Volatility\\Immutable Timeline};
\node[stage, right=0.32cm of parse] (envel) {4. Envelope \& Manifest\\Canonical Sanitization\\Manifest $\mathcal{M}$ (SHA-256)};
\node[stage, right=0.32cm of envel] (slm) {5. Air-Gapped SLM\\Quantized + KV Budget\\EBNF + Rule Check};
\node[stage, right=0.32cm of slm, fill=green!10, draw=green!60!black] (human) {6. Verified Report\\Examiner Audit\\CPP 163 / FRE 702(d)};

% Controls underneath
\node[control, below=0.25cm of evid] (c1) {ISO/IEC 27037 Custody\\SHA-256 / BLAKE3};
\node[control, below=0.25cm of block] (c2) {LED / Lei 59/2019\\Physical Air-Gap};
\node[control, below=0.25cm of parse] (c3) {ISO/IEC 27042 Analysis\\Anti-Timestomping UTC};
\node[control, below=0.25cm of envel] (c4) {Base64 Transport Layer\\Pre-Inference Strip};
\node[control, below=0.25cm of slm] (c5) {Offline \texttt{llama.cpp} / vLLM\\Dual-Agent + YARA/Sigma};
\node[control, below=0.25cm of human] (c6) {EU AI Act Annex III\\Sworn \textit{In Personam} Oath};

% Arrows
\draw[arr] (evid) -- (block);
\draw[arr] (block) -- (parse);
\draw[arr] (parse) -- (envel);
\draw[arr] (envel) -- (slm);
\draw[arr] (slm) -- (human);

\draw[dotted, thick, gray] (evid) -- (c1);
\draw[dotted, thick, gray] (block) -- (c2);
\draw[dotted, thick, gray] (parse) -- (c3);
\draw[dotted, thick, gray] (envel) -- (c4);
\draw[dotted, thick, gray] (slm) -- (c5);
\draw[dotted, thick, gray] (human) -- (c6);

\end{tikzpicture}
\caption{End-to-end air-gapped forensic ingestion pipeline reconciling local generative inference with ISO/IEC 27037/27042 custody and analysis standards, EU LED/AI Act mandates, adversarial sanitization, and courtroom verifiability.}
\label{fig:forensic_pipeline}
\end{figure*}


Fig.~\ref{fig:forensic_pipeline} consolidates this end-to-end reference pipeline, aligning every ingestion and inference stage with its mandatory cryptographic, adversarial, and legal controls.
